\documentclass[12pt,noshowpacs,nofootinbib,notitlepage,amsmath,superscriptaddress]{revtex4-1}% preprint is a loose double spaced format.
\usepackage{setspace}
\linespread{1.25}
\usepackage{hyperref}
\usepackage{amsfonts}
\usepackage{appendix}
\usepackage{mathtools}
\usepackage{ulem}
\usepackage{ifthen}
\newboolean{PrintVersion}
\setboolean{PrintVersion}{false}
\usepackage{amssymb,amstext} 
\usepackage{graphicx} 
\usepackage[caption=false]{subfig}
\usepackage{wrapfig}
\usepackage{cleveref} 
\usepackage{enumitem}
\usepackage{xcolor}
\usepackage[export]{adjustbox}%http://ctan.org/pkg/adjustbox
\usepackage{float}
\usepackage{booktabs}
\usepackage[section]{placeins}
\usepackage{siunitx}
\newcommand{\rbm}[1]{{\color{red}\bf [Robb: #1]}}
\newcommand{\mg}[1]{{\color{blue}\bf [Michael: #1]}}

\newcommand{\yb}[1]{{\color{orange}\bf [Yann: #1]}}
\newcommand{\tcb}{\textcolor{blue}}
\newcommand{\tcr}{\textcolor{red}}
\newcommand{\tco}{\textcolor{orange}}
\newcommand{\tcp}{\textcolor{purple}}

\newcommand{\package}[1]{\textbf{#1}}
\newcommand{\cmmd}[1]{\textbackslash\texttt{#1}}

\ifthenelse{\boolean{PrintVersion}}{   % for improved print quality, change some hyperref options
\hypersetup{	% override some previously defined hyperref options
%    colorlinks,%
    citecolor=black,%
    filecolor=black,%
    linkcolor=black,%
    urlcolor=black}
}{} % end of ifthenelse (no else)

\setlength{\marginparwidth}{0pt}
\setlength{\marginparsep}{0pt}
\setlength{\evensidemargin}{0.125in}
\setlength{\oddsidemargin}{0.125in}
\setlength{\textwidth}{6.375in} 
\raggedbottom
\setlength{\parskip}{\medskipamount}
\renewcommand{\baselinestretch}{1}


\begin{document}
\title{A Comment on the Stability of Charged Interacting Quark Stars}
\author{Yann Bissessur}
\email{yann.bissessur@ens-paris-saclay.fr}
\affiliation{Université Paris-Saclay, ENS Paris-Saclay, DER de Physique, 91190, 
Gif-sur-Yvette, France}
\affiliation{Department of Physics and Astronomy, University of Waterloo, Waterloo, Ontario, N2L 3G1, Canada}
\author{Michael Gammon}
\email{gammon.michael10@gmail.com}
\affiliation{Department of Physics and Astronomy, University of Waterloo, Waterloo, Ontario, N2L 3G1, Canada}
\author{Robert B. Mann}
\email{rbmann@uwaterloo.ca}
\affiliation{Department of Physics and Astronomy, University of Waterloo, Waterloo, Ontario, N2L 3G1, Canada}


%\date{\today}

\begin{abstract}
We comment on an error in the charged quark star literature \cite{Goncalves:2020joq} that resulted in erroneous solutions to the radial stability equations for charged interacting quark stars in \cite{zhang:2021}.  In this comment article we recalculate the fundamental radial eigenfrequencies for charged, interacting quark stars in general relativity and discuss any departure from previous results  \cite{zhang:2021}. Most of the qualitative trends persist, but the magnitudes of the separation between the maximum mass and stability points on a given branch can change significantly. We also note that all branches in \cite{zhang:2021} that contained no stable solutions now have stable solutions in their corrected counterparts. The net effect is that in  all of the studied cases charged quark stars have a greater range of stability than originally  predicted. We   also compare four recent observational constraints 
not included in previous studies, and find that these constraints are satisfied by a subset of the mass-radius relations we obtain.

\end{abstract}
\maketitle
\newpage

\section{Introduction}\label{sec:intro}

Thanks to recent advances in observational astrophysics and strong interaction theory, interest in studying compact stars as a tool for understanding high density matter has been rejuvenated \cite{Miller_2019,salmi2024,Horvath_2023,Abbott_2020,Pasechnik:2016wkt,Holdom:2017gdc}. Recent theoretical advances have challenged the Bodmer-Witten-Terazawa hypothesis \cite{Bodmer:1971we,Witten,Terazawa:1979hq} by showing that up-down quark matter ($ud\mathrm{QM}$) might be the ground state of baryonic matter, and thus more stable than even strange quark matter (SQM) \cite{Holdom:2017gdc}. In addition to this, several objects in tension with simple neutron equations of state coupled to general relativity (like PSR J0030+0451 \cite{Miller_2019}, PSR J0740+6620 \cite{salmi2024}, HESS J1731-347 \cite{Horvath_2023}, and GW190814 \cite{Abbott_2020}) have in particular led to a number of articles on quark stars in recent years \cite{Zhang:2020jmb,zhang:2021,gammon_2024,gammon_2025_chargedquarkstars4degb,gammon2026,banerjee_2021_quark,oikonomou_2023_colourflavour,banerjee_2021_strange,panotopoulos2022,Goncalves:2020joq,ren2020,zhang2021_3,panotopoulos2019,Jasim:2021cft}.

In 2021 Zhang, Gammon, and Mann looked at the structure and stability of quark stars with a unified interacting equation of state and a net charge \cite{zhang:2021}, presenting a host of stable charged compact objects with strong interaction effects accounted for in the equation of state \cite{Zhang:2020jmb}.  However, this paper made use of a stability equation for charged compact objects \cite{Goncalves:2020joq} that contained an error. This error, evidently a typo,  does 
 not affect the mass/radius relation for charged quark stars at all.  However the stability properties of these stars can change significantly from what was reported in \cite{zhang:2021}.

In this comment article we investigate the stability of charged, interacting quark stars, making use of the corrected stability equations. In section \ref{sec:stability} we briefly review the origin of the error  and discuss how the equations governing stability change as a result. In section \ref{sec:resultsposlam} we reproduce the results from \cite{zhang:2021} with the corrected stability data for a positive strong quark interaction constant $\lambda$. The following section correspondingly covers the negative $\lambda$ results. We also discuss the implications of the changes in the results compared to the original 
study
\cite{zhang:2021}. Section \ref{sec:discussion} summarizes qualitative trends from the preceding two sections, and finally section \ref{sec:conclusion} is used for a summary and concluding remarks. For a more thorough look at the theory governing the structure and stability of charged interacting quark stars we direct the reader to \cite{zhang:2021}.


\section{Stability of Charged Interacting Quark Stars}\label{sec:stability}

We consider a static, spherically symmetric metric ansatz
\begin{equation}
     ds^2 = -e^{2\nu(r)} dt^2 + e^{2\Lambda(r)} dr^2
  + r^2\left(d\theta^2 + \sin^2\theta\, d\phi^2\right),
  \label{metric}
\end{equation}
with the quark star interior modeled as a perfect fluid of density $\rho(r)$ and pressure $P(r)$. 
The unified, strongly interacting equation of state we use to relate these quantities is~\cite{Zhang:2020jmb}
\begin{equation}
P(r)=\frac{1}{3}(\rho(r)-4B_{\mathrm{eff}})+ \frac{4\lambda^2}{9\pi^2}\left(-1+\mathrm{sgn}(\lambda)\sqrt{1+3\pi^2 \frac{(\rho(r)-B_{\mathrm{eff}})}{\lambda^2}}\right)
\label{eos_tot}
\end{equation}
where 
\begin{equation}
\lambda=\frac{\xi_{2a} \Delta^2-\xi_{2b} m_s^2}{\sqrt{\xi_4 a_4}}
\label{lam}
\end{equation}
is the generalized strong interaction coupling constant.
Accordingly, the parameter $a_4$ represents the pQCD corrections from one-gluon exchange, $m_s$ denotes the strange quark mass, the
gap parameter $\Delta$ represents the color-superconducting gap,  and the constant $\xi$ coefficients take on particular sets of values depending on the phase of the quark matter, namely
\begin{align}
(\xi_4,\xi_{2a}, \xi_{2b}) = \left\{ \begin{array} {ll}
(( \left(\frac{1}{3}\right)^{\frac{4}{3}}+ \left(\frac{2}{3}\right)^{\frac{4}{3}})^{-3},1,0) & \textrm{2SC phase}\\
(3,1,3/4) & \textrm{2SC+s phase}\\
(3,3,3/4)&   \textrm{CFL phase}
\end{array}
\right.
\end{align}
Note that   $\lambda > 0$ provided $\Delta^2/m_s^2>\xi_{2b}/\xi_{2a}$. The parameter $B_{\rm eff}$ is the effective bag constant that accounts for the nonperturbative contribution from the QCD vacuum.  In the present manuscript we assume a typical value  $B_\mathrm{eff}=60 \; \si{MeV/fm^3}$ in order to be consistent with \cite{zhang:2021}. However, our results are also presented in a unitless form which can in principle be rescaled with any desired value of $B_\mathrm{eff}$.

 

The Tolman-Oppenheimer-Volkov (TOV) equations~\cite{Oppenheimer:1939ne,Tolman:1939jz} with charge effects included~\cite{Ray:2003gt} take the form:
\begin{align}
  \frac{\mathrm{d}q}{\mathrm{d}r} & {} = 4\pi r^2 \rho_e e^{\Lambda} \, \label{eq:TOV1} \\
  \frac{\mathrm{d}m}{\mathrm{d}r} & {} = 4\pi r^2 \rho  + \frac{q}{r}\frac{\mathrm{d}q}{\mathrm{d}r} \,  \label{eq:TOV2} \\
  \frac{\mathrm{d}P}{\mathrm{d}r} & {} = -(\rho + P)\left(4\pi r P + \frac{m}{r^2} - \frac{q^2}{r^3}\right)e^{2\Lambda}  
  + \frac{q}{4\pi r^4}\frac{\mathrm{d}q}{\mathrm{d}r}   \label{eq:TOV3} \\
   \frac{\mathrm{d}\nu}{\mathrm{d}r}  & {} = -\frac{1}{\rho + P}\left(\frac{\mathrm{d}P}{\mathrm{d}r} - \frac{q}{4\pi r^4}\frac{\mathrm{d}q}{\mathrm{d}r}\right) \, \label{eq:TOV4}
  \end{align}
where $q(r)$ and $m(r)$,  respectively, represent the charge and mass within radius $r$, $\rho_e(r)$ is the electric charge density at $r$, and
\begin{equation}
  \label{chargedTOV}
  e^{-2\Lambda(r)} = 1 - \frac{2m(r)}{r} + \frac{q(r)^2}{r^2}\,
\end{equation}
is the radial metric function.

The stability of gravitationally bound compact objects is usually studied by examining their fundamental mode eigenfrequencies. This is done by assuming a fluid element is perturbatively  displaced from its equilibrium position $r$ to $r + \delta r$, and that the perturbation has a harmonic time dependence (namely $e^{i\omega t}$). Under spherical symmetry the equations describing the perturbation are of the Sturm-Liouville form \cite{Brillante:2014lwa}, namely:
\begin{equation}\label{SLform}
\frac{d}{d r}\left[\mathcal{P} \frac{d u}{d r}\right]+\left[\mathcal{Q}+\omega^2 \mathcal{W}\right] u=0,
\end{equation}
or equivalently \cite{zhang:2021}:
\begin{equation}
    \begin{aligned}
        \frac{d u}{dr}&=\frac{\eta}{\mathcal{P}}, \\
        \frac{d \eta}{d r}&=-\left[\mathcal{Q}+\omega^2 \mathcal{W}\right] u,
    \end{aligned}
\end{equation}
where $u(r)=r^2e^{-\nu(r)} \delta r$ is the normalized displacement function. %and $e^{2\nu(r)}$ is the $tt$ metric function of the static, spherically symmetric spacetime. 

The functions $\mathcal{P}$, $\mathcal{Q}$, $\mathcal{W}$ are 
\cite{Brillante:2014lwa} 
\begin{equation}
\begin{aligned}
\mathcal{P} & =e^{\Lambda+3 \nu} r^{-2} \gamma P \\
\mathcal{Q} & =\left(\rho+P\right)\Big( r^{-2} \left(\nu^\prime\right)^2 e^{\Lambda+3 \nu}-8 \pi r^{-2} e^{3 \nu+3 \Lambda} P -r^{-6} e^{3 \nu+3 \Lambda} q^2\Big) -4 r^{-3} P^\prime e^{\Lambda+3 \nu} \\
\mathcal{W} & =e^{3 \Lambda+\nu} r^{-2}\left(\rho+P\right)
\end{aligned}
\end{equation}
where $\gamma=(1+\frac{\rho}{P})\frac{dP}{d\rho}$ is the adiabatic index.  
%and $P/\rho$ are the pressure and density, respectively. 
Making use of equation  \eqref{eq:TOV3} 
%and \eqref{eq:TOV4}, 
we arrive at a simplified final form for the $\mathcal{Q}$ equation: 
\begin{align}\label{Qcorr}
    \mathcal{Q} &= (\rho + P)r^{-2}e^{\Lambda + 3\nu}\Bigg[\nu^\prime\left(\nu^\prime + 4r^{-1}\right)- (8\pi P  + r^{-4}q^2)e^{2\Lambda}\Bigg] - \frac{q q^\prime}{\pi r^7} e^{\Lambda+3\nu}.
\end{align}

This is notably different than equation (14) in \cite{Goncalves:2020joq} which reads:
\begin{equation}\label{eq:typo}
    \mathcal{Q}=(\rho+P) r^{-2} e^{\Lambda+3 \nu}\left[\nu^\prime\left(\nu^\prime-4 r^{-1}\right)-\left(8 \pi P+r^{-4} q^2\right) e^{2 \Lambda}\right].
\end{equation}
Written this way, $\mathcal{Q}$ is missing the $-\frac{q q'}{ \pi r^7}e^{\Lambda+3\nu}$ term and has a flipped sign on the term linear in $\nu'$. The  earlier stability analysis \cite{zhang:2021}
employed \eqref{eq:typo}
but with the correct (plus) sign for the $\nu'$ term. 
%\mg{this is true for my mathematica files from the chen project, but i do wonder if his code had both typos or not.. presumably not since they must have agreed..}

From here, as before, radial stability of these stars can be determined by the sign of the fundamental mode eigenfrequency $\omega_0^2$ - radial stability is achieved when  $\omega_0^2\geq 0$. For uncharged compact stars, the point at which the fundamental radial mode eigenfrequency switches signs is always coincident with the point where $\partial M/\partial \rho_\mathrm{c}$ changes sign (here $\rho_\mathrm{c}$ is the central mass density of a star, a boundary condition). The result of this is the zero fundamental eigenfrequency point coinciding with the curve's maximum mass point.

For charged compact stars this coincidence is no longer guaranteed. The specifics of this change (whether the stability transition point moves to higher/lower central densities, and by how much) depend on the model chosen for the distribution of charge, and the strength of the strong interaction parameter $\lambda$. We use the following section to share these specifics.

 For reference we have also included four recent observational constraints in our $M$/$R$ diagrams (not present in \cite{zhang:2021}), corresponding to the interesting astrophysical observations discussed in section \ref{sec:intro}. These are objects that are in some tension with common neutron star EOSs coupled to general relativity, and thus they are interesting in the context of exotic compact objects. More specifically, in \cref{fig:modelAresults,fig:modelBresults,fig:modelCresults,fig:mA_negLam,fig:mB_negLam,fig:mC_negLam} the coloured boxes are $1\sigma$ observational estimates of mass/radii for PSR J0030+0451 (violet) \cite{Miller_2019}, PSR J0740+6620 (yellow) \cite{salmi2024}, HESS J1731-347 (gray) \cite{Horvath_2023}, and GW190814 (green) \cite{Abbott_2020}. It is intriguing that a subset of the observed mass-radius constraints for all four objects are consistent with the curves we obtain for all three charge models.

% \subsection{old}
% Thus, according to equation (14) in \cite{Goncalves:2020joq} (and hence equation (26) in \cite{zhang:2021}):
% \begin{equation}\label{eq:typo}
% \begin{aligned}
% & \mathcal{P}=e^{\Lambda+3 \nu} r^{-2} \gamma P, \\
% & \mathcal{Q}=(\rho+P) r^{-2} e^{\Lambda+3 \nu}\left[\nu^{\prime}\left(\nu^{\prime}-4 r^{-1}\right)-\left(8 \pi P+r^{-4} q^2\right) e^{2 \Lambda}\right] \\
% & \mathcal{W}=e^{3 \Lambda+\nu} r^{-2}(\rho+P),
% \end{aligned}
% \end{equation}
% where $\gamma=(1+\frac{\rho}{P})\frac{dP}{d\rho}$ is the adiabatic index, and $P/\rho$ are the pressure and density, respectively. However, written this way $\mathcal{Q}$ is missing a $-\frac{q q'}{4 \pi r^5}$ term inside the square brackets and has a flipped sign on the $\nu'$ term. It is worth mentioning that while \cite{zhang:2021} as printed has equation \eqref{eq:typo} exactly, in the numerics the errant minus sign was not present, and so only the $q'$ term is modified in the actual results. \mg{this is true for my mathematica files from the chen project, but i do wonder if his code had both typos or not.. presumably not since they must have agreed..}% \mg{there is also a sign flip so i guess there's actually more than the one term? can you comment on this yann? both \cite{Goncalves:2020joq,zhang:2021} have the minus sign in the nu' term so i just want to sure the plus sign in ours isn't a typo} 
% % \yb{If I take A.3 from [A. Brillante and I. N. Mishustin 2014 EPL 105 39001], I rename $\Phi_0\to\nu$ and substitute $p'\to TOV$, then \[\mathcal{Q}=(\rho+P) r^{-2} e^{\Lambda+3 \nu}\left[\nu^{\prime}\left(\nu^{\prime}\right)-4 r^{-1}\frac{P'}{\rho+P}-\left(8 \pi P+r^{-4} q^2\right) e^{2 \Lambda}\right] \]
% % \[=(\rho+P) r^{-2} e^{\Lambda+3 \nu}\left[\nu^{\prime}\left(\nu^{\prime}\right)-4 r^{-1}\left(-\nu'+\frac{qq'/(4\pi r^4)}{\rho+P}\right)-\left(8 \pi P+r^{-4} q^2\right) e^{2 \Lambda}\right] \]}
% This seems to be partially the result of \cite{Goncalves:2020joq} substituting 
% \begin{equation}
% \frac{dP}{d r}=-(\rho+P)\left(4 \pi r P+\frac{m}{r^2}-\frac{q^2}{r^3}\right) e^{2 \Lambda}
% \end{equation}
% into equation (30) of \cite{Brillante:2014lwa}, which is missing the $\frac{d q}{d r}$ term present in the actual TOV equation on pressure:
% \begin{equation}\label{eq:dpdrcorrect}
% \frac{dP}{d r}=-(\rho+P)\left(4 \pi r P+\frac{m}{r^2}-\frac{q^2}{r^3}\right) e^{2 \Lambda}+\frac{q}{4 \pi r^4} \frac{d q}{d r}.
% \end{equation}
% Inserting equation \eqref{eq:dpdrcorrect} into equation (30) in \cite{Brillante:2014lwa} one finds the corrected $\mathcal{Q}$ equation: 
% % \mg{qq' term has inconsistent dimension so something is wrong}
% % \begin{equation}
% %     \mathcal{Q} = (\rho + P)r^{-2}e^{\Lambda + 3\nu}\left[\nu'\left(\nu' + 4r^{-1}\right) - (8\pi P  + r^{-4}q^2)e^{2\Lambda} - \frac{qq'}{4\pi r^5}  \right].
% % \end{equation}
% % \yb{and as a result:
% \begin{equation}
%     \mathcal{Q} = (\rho + P)r^{-2}e^{\Lambda + 3\nu}\left[\nu'\left(\nu' + 4r^{-1}\right) - (8\pi P  + r^{-4}q^2)e^{2\Lambda}\right] - \frac{qq'}{4\pi r^7}e^{\Lambda+3\nu}  .
% \end{equation}
% From here, as before, radial stability of these stars can be determined by the sign of the fundamental mode eigenfrequency $\omega_0^2$ - radial stability is achieved when  $\omega_0^2\geq 0$. For uncharged compact stars the point at which the fundamental radial mode eigenfrequency switches signs is always coincident with the point where $\partial M/\partial \rho_\mathrm{c}$ changes sign (here $\rho_\mathrm{c}$ is central mass density, a boundary condition). The result of this is the zero fundamental eigenfrequency point coincides with the curve's maximum mass point.

% For charged compact stars, this coincidence is no longer guaranteed. The specifics of this change (transition point to higher/lower central densities, and by how much) depend on the model chosen for the distribution of charge (and the strength of the strong interaction parameter $\lambda$). We use the following section to share these specifics.

\section{Results for Positive $\lambda$}\label{sec:resultsposlam}
Here we present the results describing the stellar structure and the radial stability for charged, interacting quark stars with a positive strong interaction constant $\lambda$.
 %, i.e., when $\Delta^2/m_s^2>\xi_{2b}/\xi_{2a}.$ 
 In all cases the mass/radius results are unchanged (matching those in \cite{zhang:2021});  because of this we primarily comment on the location of the `stability point' (the point at which $\omega_0^2$  changes signs) relative to the maximum mass point along the mass-radius curve.
 
 \subsection{Charge Model A}\label{sec:modelAplus}
 Charge model A assumes a proportionality
 $$
 \rho_e=\alpha\rho
 $$
 between mass density and charge density \cite{zhang:2021,Goncalves:2020joq,Ray:2003gt,panotopoulos2019} (or, in unitless form, $\bar{\rho}_\mathrm{e}=\alpha \bar{\rho}$ - see appendix \ref{sec:appendix}). These results are shown along with the associated $M$/$R$ and $M/\rho_\mathrm{c}$ curves in figure \ref{fig:modelAresults}. 

 

 
 As before, the stability point coincides with the maximum mass point for the uncharged solution, and as before when charge is introduced the stability point separates from the maximum mass point and moves to a star with smaller central density. 
We find, however, that the relative size of the gap  is much smaller than previously reported
\cite{zhang:2021}. This means that more of these stars are stable than we previously thought. The effect of increasing $\lambda$ is a suppression of the separation induced by the charge, both in \cite{zhang:2021} and in figure \ref{fig:modelAresults}. 
 
Furthermore, contrary to the previous study \cite{zhang:2021}, we also see a turnaround behaviour in the stability point/maximum mass separation. By the time we get to very large charge ($\alpha=0.999$) the stability point has returned to coincidence with the maximum mass point, regardless of $\lambda$ - this is qualitatively different behaviour than what was predicted in \cite{zhang:2021}, where the separation continued to grow with charge all the way to $\alpha=0.999$. Once again we note that there are more stable stars in the parameter space than was previously thought.
 
% \subsection{Charge Model A Yann}

% Plots of mass in terms of radius and $\bar{\rho}_c$ of interacting quark stars with model-A-type charge configurations are shown in Fig.~\ref{fig:modelAresults}, where the red dots and black squares, respectively, indicate the points of maximum mass  and   zero eigenfrequency of the fundamental radial oscillation mode  for each case.  We plot results for four different values of $\alpha$.\\
% {\color{red}\clarnote{Model~A is unchanged by the corrected (self-adjoint) radial-stability analysis. The statements in this subsection remain valid; for example, for $\bar{\lambda}=0,\ \alpha=0.7$ the maximum mass occurs at $\bar{\rho}_c=3.95$ and the $\bar{\omega}_0^2=0$ point at $\bar{\rho}_c=3.28$, confirming that the marginal point lies at a \emph{smaller} central density than the maximum-mass point.}}\\
% {\color{red}\cnote{ However, for large $\alpha$, the stable mass and max mass converge again. It therefore must be possible to reach larger masses than computed by the previous study. }}
\begin{figure}[h]
 \centering
 \includegraphics[width=8 cm]{caseA_MRplot.pdf}  
\includegraphics[width=8 cm]{caseA_MRplot_largeA.pdf}  
   \includegraphics[width=8 cm]{caseA_MrhoPlot.pdf}  
  \includegraphics[width=8 cm]{caseA_MrhoPlot_largeQ.pdf}  
\caption{Plots of $\bar{M}$ vs. $\bar{R}$ (upper graphs) and $\bar{M}$ vs. $\bar{\rho}_c$ (lower graphs) of charged interacting quark stars in model A for (left) $\alpha=(0, 0.3, 0.7)$ and (right) $\alpha=0.999$. The right and top axes in each plot are the corresponding dimensional results with $B_{\rm eff}=60\, \si{MeV/fm^3}$ for illustration and consistency with \cite{zhang:2021}.
The  black, blue, violet, and red curves respectively denote   $\bar{\lambda}=(0, 0.5 , 10, \infty)$ with the sign of $\lambda$ being positive. Darker shades correspond to increasing values of $\alpha$.
The solid dots denote the maximum mass configurations, with the filled squares denoting the point at which $\bar{\omega}^2_0=0$.  Finally, the coloured boxes are $1\sigma$ observational estimates of mass/radii for PSR J0030+0451 (violet) \cite{Miller_2019}, PSR J0740+6620 (yellow) \cite{salmi2024}, HESS J1731-347 (gray) \cite{Horvath_2023}, and GW190814 (green) \cite{Abbott_2020}.} 

\label{fig:modelAresults}
\end{figure}



 \subsection{Charge Model B}
 
Charge model B directly assumes that total charge is proportional to spatial volume \cite{zhang:2021,Goncalves:2020joq,Arbanil:2015uoa} via 
$$
q(r)=Q\left(\frac{r}{R}\right)^3\equiv\beta r^3
$$ 
(or in unitless form $\bar{q}(\bar{r})=\bar{\beta}\bar{r}^3$ - see appendix \ref{sec:appendix}). The $M$/$R$ and $M$/$\rho_\mathrm{c}$ results for this charge model are shown in figure \ref{fig:modelBresults}. 

Commensurate with the previous study \cite{zhang:2021}, charge model B offsets the stability point in the opposite direction relative to charge model A (namely to larger central densities) as a nonzero charge parameter is introduced. Here this separation continues growing as $\beta$ increases - we do not observe the turnaround behaviour that was present for model A, meaning that the largest separation between maximum mass and the stability point is seen at the largest $\beta$ considered. 

Furthermore, the previous study
\cite{zhang:2021} reported that the  size of the central density separation  gets rapidly larger, then slowly becomes smaller as $\bar{\lambda}$ increases. We also observe similar behaviour here - though difficult to tell by eye,   in figure \ref{fig:modelBresults} the central density separation also follows this trend (namely growing from $\bar{\lambda}=0\to\bar{\lambda}=0.5$, and then shrinking for larger $\lambda$). In general we also see a larger average separation size,   which means the sequence contains more stable stars than reported in \cite{zhang:2021}. We also note the  truncations of the curves in the upper right panel, which  appear because of the exotic pressure profiles associated with charge model B (ie. the pressure first increases above its central value $p_\mathrm{c}$ before decreasing to 0 at finite radius). The net effect is that the $M$/$R$ curves   do not start from the origin. This behaviour was also noted in the previous study \cite{zhang:2021} where a more thorough discussion   is included.


% \tcr{ 
% Contrary to this result, we see the density separation is 
% quite insensitive to the value of $\bar{\lambda}$. This means that} relatively fewer stars are stable than what was originally reported in \cite{zhang:2021}\mg{i think this is wrong, i looked closer and it seems like $\lambda=0.5$ has larger separation than 0. it's tough to tell by eye how the separations differ between the papers because of the x-axis scale changing . can someone double check this?}.  
% \yb{the density separation is essentially due to the flatness of the $M-\rho_c$ curve near the maximum. If we compare the $\bar{\lambda}=0$ and $\bar{\lambda}=0.5$ by equal $\alpha$, the $\rho_c$ separation between max and stab is always larger for $\bar{\lambda}=0$}\mg{maybe we need to discuss briefly because in figure 2 the blue curve has a larger density separation than the black one if you measure only horizontally. i don't have the actual numbers from the code though, am i missing something?}
% \yb{yes, you are right, I was not careful enough and looked only at model A. So this should be true only for model A.}
Finally,  the curves for largest $\beta$   were originally found to have   no stability point for large enough $\bar{\lambda}$  \cite{zhang:2021} -- it was reported that the fundamental mode eigenfrequency was negative for all solutions along such curves. As we see in figure \ref{fig:modelBresults} this is no longer the case  -- all branches have at least some stable solutions, including those branches that don't start from the origin.



% \yb{This is an important feature, since the distinction with/without metastable point area
% could lead to different observed masses. Any observation of a very near extremal star
% should be interpreted model B, and rule out a whole class of models with a stable mass
% loss instead of a gain. Our models A and B are more relevant, because they really split
% into very different physical observations}

% \mg{Yann, what did you mean by `extremal star' in the above? model A is capable of supporting larger masses than model B so i don't understand what you mean. the qualitative curve shapes are very distinguishable but not on the level of a single observation.}

% \yb{If you observe a star near enough from the mass maximum (it requires knowing a way to determine it's sequence though), then it would suggest the charge profile be model B because model A would not be stable in this region. However, I am not certain there is a way to determine independently the sequence of a star, since it is specified by a charge parameter requiring to assume a model before that. I don't know a lot about astronomical observations, so maybe this is not a relevant point.}
% \mg{okay i understand what you mean now thanks}

\begin{figure}[h]
 \centering
 \includegraphics[width=8 cm]{caseB_MRplot.pdf}  
  \includegraphics[width=8 cm]{caseB_MRplot_LargeQ.pdf}  
   \includegraphics[width=8 cm]{caseB_MrhoPlot.pdf}  
  \includegraphics[width=8 cm]{caseB_MrhoPlot_largeQ.pdf}  
\caption{Plots of $\bar{M}$ vs. $\bar{R}$ (upper graphs) and $\bar{M}$ vs. $\bar{\rho}_c$ (lower graphs) of charged interacting quark stars in model B for (left) $\bar{\beta}=(0,1.5, 2.5)$   and (right) $\bar{\beta}=3.5$. The right and top axes in each plot are the corresponding dimensional results with $B_{\rm eff}=60\, \si{MeV/fm^3}$ for illustration and consistency with \cite{zhang:2021}.  The  black, blue, violet, and  red curves respectively denote  $\bar{\lambda}=(0, 0.5 , 10, \infty)$ with the sign of $\lambda$ being positive. Darker shades correspond to increasing values of $\bar{\beta}$.
The solid dots denote the maximum mass configurations, with the filled squares denoting the point at which $\bar{\omega}^2_0=0$. Finally, the coloured boxes are $1\sigma$ observational estimates of mass/radii for PSR J0030+0451 (violet) \cite{Miller_2019}, PSR J0740+6620 (yellow) \cite{salmi2024}, HESS J1731-347 (gray) \cite{Horvath_2023}, and GW190814 (green) \cite{Abbott_2020}.
}
   \label{fig:modelBresults}
\end{figure}

% \subsection{Charge Model B Yann}

% We depict the mass/radii curves for   interacting quark stars with model-B-type charge  configurations   in Fig.~\ref{fig:modelBresults}. As before, both increasing $\bar{\lambda}$ and increasing charge enlarge the mass and radius of the star, and lead to a stiffer slope for the $\bar{M}$-$\bar{\rho}_c$ curves  resulting in a smaller central density at the maximum mass point.   However we see here that  the stability point $\bar{\omega}^2_0=0$ now occurs for  larger values of 
% $\bar{\rho}_c$ than the maximum mass point, with the sizes of  the mass separation and the density separation between these two points both becoming larger with increasing $\bar{\beta}$.  Interestingly, we find the density separation size first gets  rapidly larger, then slowly becomes   smaller   as $\bar{\lambda}$ increases. For sufficiently large $\bar{\beta}$ and $\bar{\lambda}$ there \color{red} are still\color{black}are no stability points, and the configurations are radially unstable for all central densities, as shown in the upper two curves in the right diagrams of Fig.~\ref{fig:modelBresults}.\\
% \\
% {\color{red}\cnote{A corrected (self-adjoint) treatment shows that even for the largest sampled $(\bar{\beta},\bar{\lambda})$---the purple ($\bar{\lambda}=10$) and red ($\bar{\lambda}=\infty$) curves at $\bar{\beta}=3.5$ in the right diagrams of Fig.~\ref{fig:modelBresults}---a well-defined $\bar{\omega}_0^2=0$ point \emph{does} exist, located at a central density \emph{larger} than that of the maximum-mass point. Hence a stable branch is present for \emph{every} sampled $\bar{\beta}$ and $\bar{\lambda}$, and the maximum-mass configuration itself is stable. Corrected values at $\bar{\beta}=3.5$ (dimensional restitution $M[\Msun]=38.0\,\bar{M}$, $\rho_c[\mathrm{MeV/fm^3}]=240\,\bar{\rho}_c$ for $B_{\rm eff}=60\,\mathrm{MeV/fm^3}$):}}
% {\color{red}
% \begin{center}\small
% \begin{tabular}{cccc}
% \toprule
% $\bar{\lambda}$ ($\bar{\beta}=3.5$) & max mass $\bar{M}$ ($\Msun$) @ $\bar{\rho}_c$ & $\bar{\omega}_0^2=0$: $\bar{M}$ ($\Msun$) @ $\bar{\rho}_c$ & ordering \\
% \hline
% $0$      & $0.059$ ($2.3$) @ $3.26$ & $0.058$ ($2.2$) @ $4.59$ & $\bar{\omega}_0^2=0$ beyond the peak \\
% $0.5$    & $0.105$ ($4.0$) @ $1.10$ & $0.094$ ($3.6$) @ $2.59$ & $\bar{\omega}_0^2=0$ beyond the peak \\
% $10$     & $0.246$ ($9.4$) @ $0.53$ & $0.191$ ($7.3$) @ $1.14$ & \textbf{stable} (not ``unstable everywhere'') \\
% $\infty$ & $0.276$ ($10.5$) @ $0.50$ & $0.223$ ($8.5$) @ $0.93$ & \textbf{stable} (not ``unstable everywhere'') \\
% \hline
% \end{tabular}
% \end{center}}
% {\color{red}\cnote{This is an important feature, since the distinction with/without metastable point area could lead to different observed masses. Any observation of a very near extremal star should be interpreted model B, and rule out a whole class of models with a stable mass loss instead of a gain. Our models A and B are more relevant, because they really split into very different physical observations}}\\

% Another interesting feature for large $\bar{\beta}$ and large $\bar{\lambda}$ configurations  is that the stellar structure with nearly zero central density has non-zero radii and masses (i.e., $\bar{M}$-$\bar{R}$ curves do not start from the origin). This can be seen from the pressure profile plots shown in Fig.~\ref{rescaledPr_caseB}. In contrast to normal cases where $\bar{P}(\bar{r})$ decreases over increasing $\bar{r}$ monotonically as shown in Fig.~\ref{rescaledPr_caseB}a, we see in Fig.~\ref{rescaledPr_caseB}b that for some large $(\bar{\lambda}, \bar{\beta})$ sets, $\bar{P}(\bar{r})$ grows first before it decreases over $r$, resulting in a finite radius even at a tiny central density. This is because the mass density or the corresponding pressure profile has to grow (over $\bar{r}$) first to counterbalance the effect of a rapidly growing charge density when $\bar{\beta}$ is large enough.
% \begin{figure}[h]
%  \centering
%  \includegraphics[width=8cm]{CaseB_Pr1.pdf}  
%   \includegraphics[width=8cm]{CaseB_Pr2.pdf}  
% \caption{$\bar{P}$-$\bar{r}$ plots of charged interacting quark stars in model B for (a) $(\bar{\lambda}, \bar{\beta})=(0.5,2.5)$ and (b) $(\bar{\lambda}, \bar{\beta})=(0.5, 3.5), (10, 2.5)$ with a tiny rescaled central density $\bar{P}=0.00001671$.}
%    \label{rescaledPr_caseB}
% \end{figure}
% Therefore, the critical value of $\beta$ for which such a transition at near-zero central density takes place is a result of the competition between  gravitational and electrostatic forces at the origin. This can be described as%~\footnote{Note that the units conversion factor  $\zeta=1$ if $\beta$ and $\rho$ are in the same units (such as the dimensionless $\bar{\beta}$ and $\bar{\rho}$), and $\zeta\approx\bf{0.31}$ if $\beta_c$ is in units of $M_{\odot}/{\rm km}^3$ and $\rho$ is in units of $1/M^2_{\odot}$. } 
% \be
% M_{r\to0}=Q_{r\to0} \Rightarrow \rho_0\frac{4\pi}{3}r^3 =\beta_c r^3\Rightarrow \beta_c=\frac{4\pi}{3}\rho_0
% \ee
% where $\rho_0=\rho(r=0)$. Rescaling this into dimensionless form using~(\ref{rescaling_prho}) and  (\ref{rescaling_beta})
% and inserting it into the rescaled EOS Eq.~(\ref{eos_rho}), we obtain 
% %or $\zeta=\bf{??}$ if  $\beta_c$ is in units of $M_{\odot}/{\rm km}^3$ and $\rho$ is in units of $\rm MeV/fm^3$
% \be
% \bar{\beta}_c={ }\frac{4\pi}{3} \left(3\bar{P}_0+1- \frac{4}{\pi^2}\bar{\lambda} \left(-1+\sqrt{1+\frac{\pi^2}{\bar{\lambda}} {(\bar{P}_0+\frac{1}{4})}}\right) \right).
% \label{betac}
% \ee
%  As the central density $\bar{P}_0\to 0$ we have 
% \be
% \bar{\beta}_{c0}= \frac{4\pi}{3}\left(1- \frac{1}{1+\sqrt{1+\pi^2/(4\bar{\lambda})}}\right)
% \label{betac0}
% \ee
% and we see that  larger $\bar{\lambda}$  results in a smaller $\bar{\beta}_{c0}$, with the smallest  value $\bar{\beta}_{c0}= 2\pi/3\approx 2.094$ achieved for $\bar{\lambda}\to \infty$, and the largest value $\bar{\beta}_{c0}=4\pi/3\approx 4.189$   achieved for $\bar{\lambda}=0$.
%  This explains the finite $(\bar{M}, \bar{R})$ values at zero central density for large $\bar{\beta}$ and large $\bar{\lambda}$ shown in Fig.~\ref{fig:modelBresults} and Fig.~\ref{rescaledPr_caseB}.
% {\color{red}\cnote{sentence to remove for the same reason as above: the purple and red curves in the right column of Fig.~\ref{fig:modelBresults} are \emph{not} unstable for all central densities; they possess a $\bar{\omega}_0^2=0$ point at a central density larger than the maximum-mass point, and therefore a stable branch (see the corrected table above and \texttt{caseB\_MrhoPlot\_largeQ.pdf}).}}
% For very large $\bar{\beta}$ and $\bar{\lambda}$ values,  we find configurations radially \color{red}stable\color{black} unstable for all central densities, as denoted by the purple and red curves in the graphs on the right column of Fig.~\ref{fig:modelBresults}.

\subsection{Charge Model C}

In this case we fix the total charge, whose actual values ($\bar{Q}=(0, 1.538,3.076)\times 10^{-2}$, corresponding to $Q=(0,1, 2)\times 10^{20} \; \si{C}$ for a typical bag constant value $B_{\rm eff}=\rm 60\, MeV/fm^3$) are chosen to be commensurate with past work \cite{zhang:2021,Goncalves:2020joq}. The corresponding results are obtained by enumerating the charge configurations and central densities that yield the chosen fixed charge value for a given $\bar{\lambda}$. This can be implemented with interior charge distributions of type model A or model B; we choose to use the former model (in line with \cite{zhang:2021}) for the results in figure \ref{fig:modelCresults}\footnote{We find for the same fixed charges, that the results for $(\bar{M},\bar{R})$ derived from model A differ from those derived from model B on the order of $0.1\%$-$1\%$, with the former having slightly larger maximum masses.}, which is consistent with the stability points shifting to lower central densities rather than higher.  As in model B, some of the curves truncate at smaller mass/central pressure. This is due to charge model C fixing the star's total charge rather than the charge model parameter $\alpha / \beta$. Because of this, an arbitrarily small star can't necessarily support a large predetermined charge. The  curves are thus `cut off' when no solutions have the required total charge.

Here the separations are much smaller than the equivalent plot in \cite{zhang:2021} (which aligns with the results in \ref{sec:modelAplus}), which indicates that there are more stable stars  than previously reported. In the present work the fixed charges are too small to appreciably offset the stability point from the maximum mass point - the differences we do see are on the order of $<0.1\%$ for central densities, $<0.01\%$ for the radii, and $<10^{-7}$ for the unitless masses. 
%\mg{yann, is this last number also a percent?} \yb{No}

% As has mostly been the case here, the max mass/stability point separations seen in figure \ref{fig:modelCresults} are significantly smaller than the results previously presented in \cite{zhang:2021}, with the separation being hardly visible. \mg{yann can we include a zoomed in plot where the separation can be seen?} however in both cases increasing $\bar{\lambda}$ fights back against the stability gained from adding charge, moving the stability point back towards the maximum mass.

\begin{figure}[h]
 \centering
 \includegraphics[width=8 cm]{caseC_profA_MRplot.pdf}  
  \includegraphics[width=8 cm]{caseC_profA_Mrhoplot.pdf}  
\caption[Short caption]{ $\bar{M}$ vs. $\bar{R}$ (left) and $\bar{M}$ vs. $\bar{\rho}_c$ (right) plots of charged interacting quark stars in model C for $\bar{Q}=(0, 1.538,3.076)\times 10^{-2}$. The right and top axes in each plot are the corresponding dimensional results with $B_{\rm eff}=60\, \si{MeV/fm^3}$ for illustration and consistency with \cite{zhang:2021}. The  black, blue, violet, and red curves respectively denote  $\bar{\lambda}=(0, 0.5 , 10,\infty)$ with the sign of $\lambda$ being positive. Darker shades correspond to increasing values of $\bar{Q}$.
The solid dots denote the maximum mass configurations, with filled squares denoting the point at which $\bar{\omega}^2_0=0$. Finally, the coloured boxes are $1\sigma$ observational estimates of mass/radii for PSR J0030+0451 (violet) \cite{Miller_2019}, PSR J0740+6620 (yellow) \cite{salmi2024}, HESS J1731-347 (gray) \cite{Horvath_2023}, and GW190814 (green) \cite{Abbott_2020}.
}
\label{fig:modelCresults}
\end{figure}

% \subsection{Charge Model C Yann}
% In this case we fix the total charge, whose actual values are chosen to be commensurate with the previous two charge configurations.  The corresponding results are obtained by enumerating the charge configurations (like model B and central densities that yield the chosen fixed charge value for a given $\bar{\lambda}$.  We illustrate the results  in Fig.~\ref{fig:modelCresults}, where we successively chose  $\bar{Q}=(0, 1.538,3.076)\times 10^{-2}$, corresponding to $Q=(0,1, 2)\times 10^{20} C$ for a typical bag constant value $B_{\rm eff}=\rm 60\, MeV/fm^3$.  

% As before, both increasing $\bar{\lambda}$ and increasing charge enlarges the mass and radius of the star, and leads to a stiffer slope for the $\bar{M}$-$\bar{\rho}_c$ curves,  yielding a smaller central density at the maximum mass point. Interestingly, similar to model B, the stellar structure with nearly zero central density has non-zero radii and masses (i.e., $\bar{M}$-$\bar{R}$ curves do not start from the origin). This can be understood in the sense that, for a fixed $\bar{Q}$, a small radius requires a very large $\bar{\beta}$ since $\bar{Q}=\bar{\beta} \bar{r}^3$, yet our previous discussion of model B has shown both numerically and analytically that very large $\bar{\beta}$ cases can not have radii close to the origin.


  
% We can see that a larger charge pushes the $\bar{\omega}^2_0=0$ point to \color{red}larger \color{black}smaller central densities than that of the maximum mass point, \color{red} I did not read the studies referenced here, but I guess we should agree with 19 and disagree with 21. 21 should also be checked for validity\color{black}in contrast to an earlier investigation for charged strange quark stars in the context of the MIT bag model~\cite{Arbanil:2015uoa},
% but in accord with a more recent investigation of these objects with an interacting EOS~\cite{Goncalves:2020joq}.\\
% {\color{red}\cnote{The sign of the density separation for Model~C was reversed in the original: the corrected (self-adjoint) analysis places the $\bar{\omega}_0^2=0$ point at a \emph{larger} central density than the maximum-mass point, for \emph{all} sampled $\bar{\lambda}$ (both signs of $\lambda$). This is physically required, since Model~C is nothing but the Model~B charge profile ($q\propto r^3$) at a fixed total charge $Q$, and must therefore share Model~B's stability ordering (marginal point beyond the peak). The separation is very small (near coincidence) and shrinks further as $\bar{\lambda}$ increases---this last feature (see the following sentence) remains correct. Corrected values at $\bar{Q}=3.076\times10^{-2}$:}}
% {\color{red}
% \begin{center}\small
% \begin{tabular}{cccc}
% \toprule
% $\bar{\lambda}$ ($\bar{Q}=3.076\times10^{-2}$) & max mass @ $\bar{\rho}_c$ & $\bar{\omega}_0^2=0$ @ $\bar{\rho}_c$ & $\Delta\bar{\rho}_c$ \\
% \hline
% $0$      & $4.40$ & $4.54$ & $+0.14$ \\
% $0.5$    & $2.82$ & $2.86$ & $+0.04$ \\
% $10$     & $1.70$ & $1.70$ & $+0.005$ \\
% $\infty$ & $1.49$ & $1.49$ & $+0.002$ \\
% \hline
% \end{tabular}
% \end{center}}
% Interestingly, we observe that a larger $\bar{\lambda}$ tend to move the $\bar{\omega}^2_0=0$ point closer to the maximum mass point. This means that for a fixed total charge strong interaction effects tend to offset charge  destabilization effects on radial stability, as expected from the confining nature of the strong interaction. 
% %\tred{[Why this feature not manifest in other charge models? Perhaps show up in Mass vs $\rho_c$ plot?]}


\section{results for negative $\lambda$}\label{sec:resultsneglam}
In the previous treatment of model A with $\lambda<0$ \cite{zhang:2021}, the added charge was found to always induce a separation between the two points of interest, and as the magnitude of  (negative) $\lambda$ grew, an increase was seen in the separation between points in central density space. The same behaviour is seen here, but with much smaller density separations on average than those presented in \cite{zhang:2021} (see figure \ref{fig:mA_negLam}). As expected in model A,  the maximum mass points are  at larger central densities than the stability points (in both \cite{zhang:2021} and the present work). This density separation concurrently manifests as a decrease in separation in $M$/$R$ space in \cite{zhang:2021}, which we find is also the case here, just on smaller separation scales. This isn't inconsistent with the former result as a large change in central density doesn't necessarily correspond to a large change in mass.
  However, unlike \cite{zhang:2021}, we once again see that the effect of large charge in model A is to force the separation back to 0, 
% \mg{is what i said correct yann? or are they offset a little bit still?}\yb{they still ofset a little bit, but in an amount that should be within the uncertainty of not reaching the exact charge limit.}
even when $\lambda$ is negative and has a large magnitude. In the previous study \cite{zhang:2021} the large charge limit was found to have very low differences in central density between maximum mass and stability points, but had a large mass separation between those same points. Employing the corrected stability equations, we do not observe this.   These new predictions allow for more stable stars that are strongly charged.

For model B with $\lambda<0$, introducing charge was found to offset the stability point to higher central densities \cite{zhang:2021}  (as it does here), but only appreciably so when the magnitude of $\lambda$ is small. This means that making $\lambda$ more negative here suppresses the separation induced by charge model B. We see analogous behaviour in figure \ref{fig:mB_negLam}, but with larger average separations than those in \cite{zhang:2021}. This means that relative to \cite{zhang:2021} we predict more stable stars in the sequence.

 When $\lambda<0$, the results for charge model C in the previous investigation \cite{zhang:2021} indicated that
separations between maximum mass and stability points are made larger as the magnitude of $\lambda$ increases, in concert with the charge, which also induces separation. Together this is manifest as  charge effects being most effective when $\lambda$ has a large magnitude and is negative. In our updated analysis we find that the behaviour is similar (see figure \ref{fig:mC_negLam}), but with the separation sizes generally being much smaller (again leading to more of the stars being stable than we previously thought). 

We have again included the four recent observational constraints (not present in \cite{zhang:2021}) in all of our $M$/$R$ diagrams for the $\lambda<0$ case as well. However in this case we find that  the observed mass-radius constraints for all four objects are not always consistent with the curves we obtain for the different charge models, due to negative $\lambda$ pushing results to smaller masses and radii).


\begin{figure}[htb]
 \centering
  \includegraphics[width=8 cm]{caseA_MRplot_negLam.pdf}  
    \includegraphics[width=8 cm]{caseA_MRPlot_negLam_LQ.pdf}  
   \includegraphics[width=8 cm]{caseA_MrhoPlot_negLam.pdf}  
      \includegraphics[width=8 cm]{caseA_MrhoPlot_negLam_LQ.pdf}  
\caption{Plots of $\bar{M}$ vs. $\bar{R}$ (upper graphs) and $\bar{M}$ vs. $\bar{\rho}_c$ (lower graphs) of charged interacting quark stars in model A for (left) $\alpha=(0, 0.3, 0.7)$ and (right) $\alpha=0.999$. The right and top axes in each plot are the corresponding dimensional results with $B_{\rm eff}=60\, \si{MeV/fm^3}$ for illustration and consistency with \cite{zhang:2021}. The  black, blue, and violet curves respectively denote $\bar{\lambda}=(0, 0.5 , 10)$ with the sign of $\lambda$ being negative. Darker shades correspond to increasing values of $\alpha$. The solid dots denote the maximum mass configurations, with the filled squares representing where $\bar{\omega}^2_0=0$.  Finally, the coloured boxes are $1\sigma$ observational estimates of mass/radii for PSR J0030+0451 (violet) \cite{Miller_2019}, PSR J0740+6620 (yellow) \cite{salmi2024}, HESS J1731-347 (gray) \cite{Horvath_2023}, and GW190814 (green) \cite{Abbott_2020}. Note that for the violet curves in the $\bar{M}$ vs. $\bar{\rho}_c$ plots, we rescaled the x-axis as $\bar{\rho}_c\to\bar{\rho}_c/5$ for clarity.}
\label{fig:mA_negLam}
\end{figure}

\begin{figure}[h]
 \centering
 \includegraphics[width=8 cm]{caseB_MRplot_negLam.pdf}  
     \includegraphics[width=8 cm]{caseB_MRplot_negLam_LQ.pdf}  
   \includegraphics[width=8 cm]{caseB_MrhoPlot_negLam.pdf}  
     \includegraphics[width=8 cm]{caseB_MrhoPlot_negLam_LQ.pdf}  
\caption{Plots of $\bar{M}$ vs. $\bar{R}$ (upper graphs) and $\bar{M}$ vs. $\bar{\rho}_c$ (lower graphs) of charged interacting quark stars in model B for (left) $\bar{\beta}=(0,1.5, 2.5)$  and (right) $\bar{\beta}=3.5$. The right and top axes in each plot are the corresponding dimensional results with $B_{\rm eff}=60\, \si{MeV/fm^3}$ for illustration and consistency with \cite{zhang:2021}. The black, blue, and violet curves respectively denote $\bar{\lambda}=(0, 0.5 , 10)$ with the sign of $\lambda$ being negative. Darker shades correspond to increasing values of $\bar{\beta}$. 
The solid dots denote the maximum mass configurations, with the filled squares representing where $\bar{\omega}^2_0=0$.  Finally, the coloured boxes are $1\sigma$ observational estimates of mass/radii for PSR J0030+0451 (violet) \cite{Miller_2019}, PSR J0740+6620 (yellow) \cite{salmi2024}, HESS J1731-347 (gray) \cite{Horvath_2023}, and GW190814 (green) \cite{Abbott_2020}. Note that for the violet curves in the $\bar{M}$ vs. $\bar{\rho}_c$ plots, we rescaled the x-axis as $\bar{\rho}_c\to\bar{\rho}_c/5$ for clarity.}
\label{fig:mB_negLam}
\end{figure}

\begin{figure}[h]
 \centering
   \includegraphics[width=8 cm]{caseC_profA_MRplot_negLam.pdf}  
  \includegraphics[width=8 cm]{caseC_profA_Mrhoplot_negLam.pdf}  
\caption{$\bar{M}$ vs. $\bar{R}$ (left) and $\bar{M}$ vs. $\bar{\rho}_c$ (right) of charged interacting quark stars in model C for $\bar{Q}=(0, 1.538,3.076)\times 10^{-2}$. The right and top axes in each plot are the corresponding dimensional results with $B_{\rm eff}=60\, \si{MeV/fm^3}$ for illustration and consistency with \cite{zhang:2021}. The  black, blue, and violet curves respectively denote  $\bar{\lambda}=(0, 0.5 , 10)$ with the sign of $\lambda$ being negative. Darker shades correspond to increasing values of $\bar{Q}$. The solid dots denote the maximum mass configurations, with filled squares denoting the point at which $\bar{\omega}^2_0=0$. Finally, the coloured boxes are $1\sigma$ observational estimates of mass/radii for PSR J0030+0451 (violet) \cite{Miller_2019}, PSR J0740+6620 (yellow) \cite{salmi2024}, HESS J1731-347 (gray) \cite{Horvath_2023}, and GW190814 (green) \cite{Abbott_2020}.  Note that for the violet curves in the $\bar{M}$ vs. $\bar{\rho}_c$ plots, we rescaled the x-axis as $\bar{\rho}_c\to\bar{\rho}_c/5$ for clarity.
}
\label{fig:mC_negLam}
\end{figure}




\section{Discussion}\label{sec:discussion}

We have analyzed the
stability properties of charged quark stars using equation \eqref{SLform} with $\cal Q$ given by \eqref{Qcorr}, correcting the equation employed in a previous study  \cite{Goncalves:2020joq,zhang:2021}. The individual cases (consisting of charge model, charge strength, and size/sign of $\lambda$) in general differ in their effects on the separation between the stability point and maximum mass point, but some general trends can be observed.

First, the relative direction that the stability point moves with respect to the maximum mass point is unchanged from what was previously reported \cite{zhang:2021}. In model A (and C, which implicitly uses model A) the stability points move toward smaller central densities, whereas in model B they move to larger central densities. The size of the induced separation between the two points   in general differs from what was obtained previously in \cite{zhang:2021}. 

In general  the separation size in model A increases as the charge parameter $\alpha$ increases, but eventually turns around and closes for maximal $\alpha$. Making $\lambda$ more positive in these cases suppresses this separation, and accordingly making $\lambda$ more negative can contribute to the separation. Barring the turnaround at large charge which is unique to the present investigation, these general trends are in accordance with the earlier work \cite{zhang:2021} (however the separation sizes are smaller than previously reported in \cite{zhang:2021}).

For model B, separation size also increases as $\beta$ increases but we do not observe any kind of turnaround behaviour in the range of $\beta$ considered. The effects of $\lambda$ are slightly more complicated than the model A case -- making $\lambda$ more positive first contributes to an increased separation from $\bar{\lambda}=0\to\bar{\lambda}=0.5$, but then suppresses separation when increased further. On the other hand, making $\lambda$ more negative strictly suppresses the separation induced by $\beta$. It is worth noting that \cite{zhang:2021} also saw this turnaround behaviour for $\lambda>0$ under charge model B \cite{zhang:2021}, and negative $\lambda$ also had a suppressing effect there. The most notable difference is the fact that in this work all model B branches have stable stars, including when $\lambda\to\infty$ (which was not the case in \cite{zhang:2021}).

In general the separation sizes for the fixed charges considered in model C are very small in the present paper and are hardly perceptible on an $M$/$R$ plot, but since these solutions are implicitly using model A, we know that the qualitative effects of $\lambda$ and (small) charge parameter $\alpha$ in both our study and \cite{zhang:2021} match those already discussed for model A. The unique feature here (in contrast with \cite{zhang:2021}) is that once again all branches contain stable star solutions, even those starting relatively near the maximum mass point.

\section{Conclusion}\label{sec:conclusion}

We have numerically recalculated the stability properties of charged interacting quark stars, with the stability equations from \cite{Goncalves:2020joq} (propagating to \cite{zhang:2021}) corrected. Many of the general trends of the parameter space remain unchanged, but the actual magnitudes of the separations between the stability point and maximum mass can change significantly.

The most interesting qualitative changes to the results include ``turnaround" behaviour of the separation which reapproaches 0 for maximal model A charge parameter $\alpha$, and the fact that model B branches which were entirely unstable in \cite{zhang:2021} all contain stable star solutions in the corrected results. In all cases these changes allow for more stable stars than was reported in \cite{zhang:2021}.

 It is clear from \cref{fig:modelAresults,fig:modelBresults,fig:modelCresults,fig:mA_negLam,fig:mB_negLam,fig:mC_negLam} that, for particular choices of the strong interaction parameter (such as $\lambda > 0$ and $0\leq \bar{\lambda} \leq 1$), a subset of these solutions are stable and also pass through all four of the included observational constraints. A larger number of stable solutions satisfy only some of the constraints, and a larger amount still satisfy none. There is no guarantee that these objects actually are (charged) quark stars, however if they are, the data is consistent with an interacting equation of state for particular values of the strong interaction parameter $\lambda$.
 


\section*{Acknowledgements}
This work was supported in part by the Natural Sciences and Engineering Research Council of Canada.



\bibliographystyle{unsrt}

\cleardoublepage
\phantomsection  
\renewcommand*{\refname}{References}


% \printbibliography[title=References]
\bibliography{bib}

\appendix
\section{Units and Scaling}\label{sec:appendix}
In all cases we numerically solve the TOV equations in a unitless form where the stellar variables have been rescaled using the effective bag constant $B_\mathrm{eff}$. More specifically, if we perform the following rescalings:
\begin{equation}
    \bar{\rho}=\frac{\rho}{4\,B_{\rm eff}}, \,\, \bar{P}=\frac{P}{4\,B_{\rm eff}}, \,\,\bar{\lambda}=\frac{\lambda^2}{4B_{\rm eff}}
\end{equation}
along with
\begin{align}
\bar{m}&=m{\sqrt{4\,B_{\rm eff}}}, \quad \bar{r}={r}{\sqrt{4\,B_{\rm eff}}},\\
\bar{q}&={q}{\sqrt{4\,B_{\rm eff}}}, \quad \bar{\rho}_e=\frac{\rho_e}{4\,B_{\rm eff}}
\end{align}
the TOV equations \eqref{eq:TOV1}-\eqref{eq:TOV4} can be rewritten in unitless form simply by replacing all variables with their unitless (barred) counterparts. Regarding the rescaling of charge model parameters, $\alpha$ is unitless to begin with. However $\beta$ and $Q$ rescale as
\begin{equation}
    \bar{\beta}=\frac{\beta}{4B_{\rm eff}}, \quad \bar{Q}={Q}{\sqrt{4\,B_{\rm eff}}}.
\end{equation}


\end{document} 

\rbm{I don't get the above; instead I get
\begin{align}
    \mathcal{Q} &= (\rho + P)r^{-2}e^{\Lambda + 3\nu}\Bigg[\nu^\prime\left(\nu^\prime + 4r^{-1}\right)- (8\pi P  + r^{-4}q^2)e^{2\Lambda}\Bigg] - \frac{q q^\prime}{\pi r^7} e^{\Lambda+3\nu}
    \nonumber
\end{align}
Note the difference in the last term.
Please check, and chack against the code.
}
\yb{You are right, the substitution of the conservation equation into (11) should give your version without the $\frac{1}{4}$ factor. It must be a typo that appears only in the text, because the code uses the general formula that I derived for the NLED article and that recovers yours in linear limit.}